\documentclass[11pt]{amsart}

\usepackage[T1]{fontenc}
\usepackage{lmodern}
\usepackage{microtype}
\usepackage{amsmath,amssymb,amsthm}
\usepackage{booktabs}
\usepackage[hidelinks]{hyperref}

\hypersetup{
  pdftitle={A Unit-Cube Box-Visibility Representation of K9},
  pdfsubject={Box-visibility representations of complete graphs},
  pdfkeywords={
    box visibility,
    unit cubes,
    complete graphs,
    geometric graph representations
  }
}

\newtheorem{theorem}{Theorem}
\newtheorem{corollary}[theorem]{Corollary}

\title[A unit-cube representation of $K_9$]
{A Unit-Cube Box-Visibility Representation of $K_9$}

\date{}

\subjclass[2020]{05C62, 68U05}
\keywords{
box visibility,
unit cubes,
complete graphs,
geometric graph representations
}

\begin{document}

\begin{abstract}
We give explicit coordinates for nine pairwise disjoint axis-parallel
unit cubes whose box-visibility graph is $K_9$. This refutes
Conjecture~26 of Fekete and Meijer. Together with their theorem that
$K_{10}$ has no such representation, this determines $9$ as the
maximum order of a complete graph representable by unit cubes.
\end{abstract}

\maketitle

A box-visibility representation assigns pairwise disjoint
axis-parallel boxes in $\mathbb R^3$ to the vertices of a graph, with
two vertices adjacent exactly when their boxes can be joined by an
axis-parallel cylindrical channel of positive radius whose ends lie
in the box interiors and which meets no other box. Fekete and Meijer
constructed such a representation of $K_8$ by unit cubes, proved that
none exists for $K_{10}$, and conjectured the same for $K_9$
\cite[Fig.~5, Theorem~25, and Conjecture~26]{FM}.

\begin{theorem}\label{thm:main}
There exist nine pairwise disjoint axis-parallel unit cubes whose
box-visibility graph is $K_9$. In particular, Conjecture~26 of Fekete
and Meijer is false.
\end{theorem}

\begin{proof}
It is convenient to work at scale $16$. For $i=0,\ldots,8$, let
\[
 C_i
 =
 [x_i,x_i+16]
 \times
 [y_i,y_i+16]
 \times
 [z_i,z_i+16],
\]
where
\begin{equation}\label{eq:coordinates}
\begin{array}{c|rrrrrrrrr}
 i   &0&1&2&3&4&5&6&7&8\\ \hline
 x_i &0&9&-8&10&6&2&7&-5&5\\
 y_i &0&2&3&-1&-7&11&-6&8&1\\
 z_i &0&17&19&34&51&53&68&85&102.
\end{array}
\end{equation}
The $z$-intervals overlap only for the pairs
$\{1,2\},\{2,3\},\{4,5\},\{5,6\}$. The first two pairs are strictly
separated in the $x$-direction, and the last two in the $y$-direction.
Thus the cubes are pairwise disjoint. All minima and maxima are
pairwise distinct in each coordinate direction.

Write
\[
\begin{aligned}
 \ell_x(a,b)&=\{(t,a,b):t\in\mathbb R\},\\
 \ell_y(a,b)&=\{(a,t,b):t\in\mathbb R\},\\
 \ell_z(a,b)&=\{(a,b,t):t\in\mathbb R\}.
\end{aligned}
\]
The following table gives a visibility axis for every unordered pair;
the label $ij:d(a,b)$ means the line $\ell_d(a,b)$ for the pair
$\{C_i,C_j\}$.

\begin{center}
\footnotesize
\renewcommand{\arraystretch}{1.12}
\setlength{\tabcolsep}{7pt}
\begin{tabular}{lll}
\toprule
$01:z(19/2,5/2)$
&
$02:z(1/2,7/2)$
&
$03:z(21/2,1/2)$
\\
$04:z(13/2,1/2)$
&
$05:z(17/2,23/2)$
&
$06:z(17/2,19/2)$
\\
$07:z(17/2,21/2)$
&
$08:z(11/2,3/2)$
&
$12:x(7/2,39/2)$
\\
$13:z(21/2,5/2)$
&
$14:z(19/2,5/2)$
&
$15:z(19/2,23/2)$
\\
$16:z(19/2,19/2)$
&
$17:z(19/2,21/2)$
&
$18:z(37/2,31/2)$
\\
$23:x(7/2,69/2)$
&
$24:z(13/2,7/2)$
&
$25:z(5/2,23/2)$
\\
$26:z(15/2,19/2)$
&
$27:z(-9/2,17/2)$
&
$28:z(11/2,7/2)$
\\
$34:z(21/2,-1/2)$
&
$35:z(21/2,23/2)$
&
$36:z(21/2,19/2)$
\\
$37:z(21/2,21/2)$
&
$38:z(23/2,21/2)$
&
$45:y(13/2,107/2)$
\\
$46:z(15/2,-11/2)$
&
$47:z(13/2,17/2)$
&
$48:z(13/2,3/2)$
\\
$56:y(15/2,137/2)$
&
$57:z(5/2,23/2)$
&
$58:z(23/2,23/2)$
\\
$67:z(15/2,17/2)$
&
$68:z(15/2,3/2)$
&
$78:z(11/2,17/2)$
\\
\bottomrule
\end{tabular}
\end{center}

\newpage

For $d\in\{x,y,z\}$, write $p_d$ for the $d$-coordinate of
$p\in\mathbb R^3$, and let $\pi_d$ denote projection onto the two
coordinates transverse to $d$, ordered as in the definitions of
$\ell_d$. Fix an entry $ij:d(a,b)$ and put $q=(a,b)$. After exchanging
$i$ and $j$ if necessary, let $u$ be the upper $d$-face of the lower
cube and $v$ the lower $d$-face of the upper cube. Direct substitution
in~\eqref{eq:coordinates} gives, for every table entry,
\begin{align*}
 \operatorname{dist}_\infty\!\left(
 q,\mathbb R^2\setminus\operatorname{int}(\pi_d(C_h))\right)
 &\ge \frac12
 &&(h\in\{i,j\}),\\
 \operatorname{dist}_\infty(q,\pi_d(C_k))
 &\ge \frac12
 &&\left(
 k\notin\{i,j\},\
 I_k^d\cap[u-\tfrac14,v+\tfrac14]\ne\varnothing
 \right),
\end{align*}
where $I_k^d$ is the $d$-interval of $C_k$. These are exact finite
checks: the face coordinates are integers and the listed transverse
coordinates are half-integers. Consequently the cylinder
\[
 \Gamma_{ij}
 =
 \left\{
 p\in\mathbb R^3:
 u-\tfrac14\le p_d\le v+\tfrac14,\
 \|\pi_d(p)-q\|_2\le\tfrac14
 \right\}
\]
has its end disks in the interiors of $C_i$ and $C_j$ and is disjoint
from every other cube. Thus every pair of cubes is visible. Dividing
all coordinates by $16$ gives unit cubes and visibility channels of
radius $1/64$.
\end{proof}

The coordinates were found computationally, but the displayed
certificate and its verification are exact.

\begin{corollary}
The largest $n$ for which $K_n$ has a box-visibility representation
by axis-parallel unit cubes is $n=9$.
\end{corollary}

\begin{proof}
Theorem~\ref{thm:main} gives the lower bound. Fekete and Meijer proved
that $K_{10}$ has no such representation~\cite[Theorem~25]{FM}.
A representation of any $K_n$ with $n>10$, restricted to any ten
cubes, would give a representation of $K_{10}$, a contradiction.
\end{proof}

\begin{thebibliography}{1}

\bibitem{FM}
S.~P. Fekete and H. Meijer,
\emph{Rectangle and box visibility graphs in 3D},
Internat. J. Comput. Geom. Appl.
\textbf{9} (1999), no.~1, 1--27,
\href{https://doi.org/10.1142/S0218195999000029}
{doi:10.1142/S0218195999000029}.

\end{thebibliography}

\end{document}